\subsection{Three visible ways continuity can fail}

The pictures below separate three common phenomena. The algebraic formulas may differ, but the graph reveals the type of failure immediately.

\begin{center}
\begin{tikzpicture}[x=0.75cm,y=0.75cm]
  % removable
  \begin{scope}[xshift=0cm]
    \draw[->] (-2,0) -- (2,0) node[right] {$x$};
    \draw[->] (0,-1) -- (0,2.5) node[above] {$y$};
    \draw[thick,domain=-1.7:1.7,samples=2] plot (\x,{0.6*\x+1});
    \fill[white] (0.7,1.42) circle (2.8pt);
    \draw[thick] (0.7,1.42) circle (2.8pt);
    \fill (0.7,2.05) circle (2.2pt);
    \node at (0,-1.45) {removable};
  \end{scope}

  % jump
  \begin{scope}[xshift=5.2cm]
    \draw[->] (-2,0) -- (2,0) node[right] {$x$};
    \draw[->] (0,-1) -- (0,2.5) node[above] {$y$};
    \draw[thick] (-1.7,0.65) -- (0,0.65);
    \fill (0,0.65) circle (2.2pt);
    \draw[thick] (0,1.75) -- (1.7,1.75);
    \fill[white] (0,1.75) circle (2.8pt);
    \draw[thick] (0,1.75) circle (2.8pt);
    \node at (0,-1.45) {jump};
  \end{scope}

  % infinite
  \begin{scope}[xshift=10.4cm]
    \draw[->] (-2,0) -- (2,0) node[right] {$x$};
    \draw[->] (0,-1) -- (0,2.5) node[above] {$y$};
    \draw[dashed] (0,-0.8) -- (0,2.3);
    \draw[thick,domain=-1.7:-0.25,samples=60] plot (\x,{-0.35/\x});
    \draw[thick,domain=0.25:1.7,samples=60] plot (\x,{0.35/\x});
    \node at (0,-1.45) {infinite};
  \end{scope}
\end{tikzpicture}
\end{center}

\begin{interpretationbox}[Classify the failure geometrically]
A removable discontinuity has a common nearby height but the wrong or missing value. A jump has different one-sided limits. An infinite discontinuity occurs when values become unbounded near a vertical line.
\end{interpretationbox}
